%!TEX root = ../hutbthesis_main.tex

\chapter{结论与展望}

本文将基于人机交互强化的固定翼飞行器着陆控制系统当作研究对象，针对系统组成、需求分析、动力学模型和仿真平台建立、多工况仿真结果分析，还有人机交互强化策略设计等方面展开了研究。

研究结果表明，固定翼飞行器着陆阶段受速度、下沉率、姿态角、侧偏量以及地形环境等多种因素共同影响，是飞行全过程中控制要求较高、风险较集中的阶段。通过建立着陆控制动力学模型和仿真平台，可以较好地反映飞行器在进近、拉平和接地区域的状态变化规律。基准工况校核结果表明，所建模型具有较好的可运行性和合理性，能够满足后续工况分析需要。

多工况仿真得出的结果显示，在不同地形状况之下，飞行器着陆末段的响应呈现出显著的差别。平坦地形工况整体具备最高的稳定性，能够当作参考基线，谷地地形工况有着较好的表现，起伏地形、山脊地形会使得姿态修正频率有所增加，阶跃地形对于近地判断、接地控制产生的影响最为突出，属于最为典型的不利工况。由此能够看出，地形环境发生改变，对于固定翼飞行器着陆控制品质有着重要作用。

基于此本文形成了一种人机交互强化思路，将风险识别、状态提示、操纵辅助、约束保护整合在一起。研究显示，此策略可在一定程度上提高操纵者对于复杂工况的感知、应对能力，对于改进着陆安全性和稳定性有着积极作用。

因为研究条件、论文篇幅方面存在限制，所以本文存在一些不足之处。其一，本文主要运用仿真分析方法，对于真实飞行器复杂的气动特性、实际扰动环境的考量不够全面，后续能够结合更高保真度模型进一步优化研究。其二，本文针对起落架接地缓冲过程主要做了简化处理，后续可结合接地冲击、缓冲系统动力学展开更细致的分析。其三，本文所提出的人机交互强化策略主要在仿真平台上得到了验证，后续可结合半实物仿真或者飞行试验进一步检测其有效性。

当下，针对固定翼飞行器着陆控制系统的人机交互强化这一研究领域而言，存在着比较可观的拓展可能性。在后续的发展进程当中，可以把更为复杂的工况状况、精度更高的模型、更为完善的交互设计方式相互结合起来。通过这种结合方式，能够提高固定翼飞行器在复杂环境里的着陆安全程度，进而提高其在工程应用方面的价值。
